%---------------------------------------------------------------------
\section{Conclusion}

\subsection*{Hangboarding}

Hangboarding is an evidence-supported supplemental training method for intermediate-to-advanced climbers. The meta-analytic evidence \cite{SO23} confirms consistent short-term improvements in dynamometric finger strength, RFD, and forearm endurance (SDM 0.41--1.23), but no trial has yet demonstrated grade improvement attributable to hangboarding alone. Max-hang protocols ($\geq$80\% MFS) best develop MFS; repeater protocols (60--80\% MFS) best develop stamina; 80\% MFS offers both. Daily low-intensity hangs (Abrahangs) produced a positive but non-significant trend versus climbing-only controls ($p = 0.12$) in the sole observational cohort; RCT confirmation is required before broad recommendation. Injury risk is experience-dependent: experienced climbers may benefit from a protective effect, while rapidly advancing less-experienced climbers face elevated incidence. Critical evidence gaps remain: female-cohort data, trials $>$10 weeks with injury tracking, and any functional on-wall grade outcome.

\subsection*{Weekly Training Integration}

Session sequencing and weekly structure are predominantly BIOPLAUSIBLE or FOLKLORE\_VERIFIED (ESS 1--2); no RCT has tested a complete weekly training template as a unit. The $\geq$48 h inter-session minimum between high-intensity finger sessions ($\geq$80\% MFS) is a study-design convention, not an empirically optimised prescription. Hangboard max-strength work must precede, not follow, limit bouldering when MFS is the target. Evidence for concurrent-training interference is well-established in lower-body strength-endurance models \cite{WMR+12} and is bioplausibly transferable to climbing, but no RCT has measured its magnitude in climbers. Deload indicators (RPE elevation, rolling force decrement, HRV trends) are extrapolated from general resistance training and have not been validated in climbing populations.

\subsection*{Supplementation}

No supplement in this review has been tested with a functional climbing outcome as the primary endpoint. Within that constraint: \textbf{collagen + vitamin C} (15\,g hydrolysed collagen $+$ 48\,mg ascorbic acid 1\,h pre-exercise) produces the most direct mechanism-to-tissue link for pulley and tendon support, though evidence remains surrogate (P1NP biomarker; ESS 3--4); \textbf{beta-alanine} (3.2--6.4\,g/day $\geq$4 weeks) has the strongest climbing-specific RCT signal for endurance tasks of 60--240\,s \cite{SNWK21} (ESS 7, but $n=15$); \textbf{caffeine} (3\,mg/kg) showed null results on all proxy climbing metrics in the only published climbing RCT \cite{RLMAMR+26} (ESS 7, $n=13$); \textbf{creatine} is evidence-backed for strength and power in general populations but carries a body-mass penalty that directly offsets performance on gravity-dependent movements; \textbf{omega-3} and \textbf{magnesium} lack climbing-specific controlled evidence and remain bioplausible at most for sub-populations with documented deficiency.

\subsection*{Injury Management}

A2 pulley injuries are the most common acute climbing finger injury; the corrected Sch\"{o}ffl Grade II--IV classification, pulley-protection splint (PPS) as primary conservative device (88.4\% return at mean 8.8\,months; \cite{SS16}), and Grade IVa/IVb surgical threshold are key updates from the initial framework. Flexor tenosynovitis (second most common; favourable natural history under load reduction, \cite{MSM+22}) and PIP capsulitis/synovitis (third most common; staged algorithm with 72\% return, \cite{SLL+25}) represent previously absent coverage that substantially broadens the section's clinical utility. Lumbrical tears are identified by the pathognomonic lumbrical stress test and managed with pocket-avoidance and progressive loading (\cite{LSS+18}). Adolescent physeal injuries (the dominant injury in $<$16-year-old climbers) require immediate rest and imaging; the adult pulley framework is contraindicated (\cite{SSF+21}). Injury risk is experience-dependent: a bidirectional pattern (elevated risk in intermediate climbers, apparent protection in highly experienced climbers) is observationally supported \cite{SGJ23} (ESS 5). Critical gaps include absence of RCTs comparing PPS versus H-taping, progressive loading versus rest-only for Grade I--II, NSAID effects on pulley healing, and long-term OA progression after capsulitis episodes.

\subsection*{Recovery}

No RCT has tested any recovery modality specifically against climbing performance or pulley/tendon injury incidence. Within that constraint: \textbf{sleep} $\geq$8\,h is the recovery modality with the strongest evidence base across sports \cite{MSB+14}; \textbf{cold water immersion} is evidence-backed for DOMS reduction but should be avoided immediately post-strength/hangboard sessions as it attenuates anabolic signalling \cite{RRM+15}; \textbf{massage and foam rolling} reduce perceived soreness (SMD $-$0.72; \cite{DDT+18}) but do not address tendon recovery on the collagen-remodelling timeline; \textbf{active recovery} (20--30\,min aerobic at RPE 3--4) accelerates lactate clearance and has one climbing-specific controlled study (Mourot et al.\ 2015; \toverify: DOI unresolved) with a functional outcome. The threshold between active recovery and additional connective tissue stress is not empirically defined for any climbing grip position.

\subsection*{Evidence Gaps and Pending Sections}

Across all domains, the critical shared gaps are: absence of female-cohort climbing data; lack of long-duration ($>$10 week) trials with concurrent injury tracking; and no trial measuring grade improvement as a primary outcome. Performance predictors (finger strength, pull-up metrics, and their correlation with climbing grade), equipment ROI, scientific measurement methodology (Tindeq Critical Force, W', force asymmetry), and skincare are addressed in forthcoming sections.

Until climbing-specific high-quality RCTs exist for these domains, practitioner folklore should be applied as biologically plausible guidance calibrated to individual tolerance, not as protocol with established dosing certainty.
